% Heya mate, this is not finished yet. Just so you know... Oh, I almost forgot, this document is provided AS IS without any warranty of any sort. The content might be total bullshit and should be used with caution. So if it burns down your house, kills your cat or opens a portal to another dimension, you are at fault. Don't come complaining I didn't warn ya. With that out of the way, enjoy the document!

\documentclass[12pt,a4paper,ngerman]{scrartcl}
\usepackage[utf8]{inputenc} % üäö directly in source
\usepackage[ngerman]{babel}
\usepackage{float}
\usepackage{pdfpages}
\usepackage{titlesec}
\usepackage[colorlinks, linkcolor = black, citecolor = black, filecolor = black, urlcolor = blue]{hyperref} 

%Polynomdivision
\usepackage{polynom}


%4 Ebenen
\setcounter{secnumdepth}{4}

\titleformat{\paragraph}
{\normalfont\normalsize\bfseries}{\theparagraph}{1em}{}
\titlespacing*{\paragraph}
{0pt}{3.25ex plus 1ex minus .2ex}{1.5ex plus .2ex}


%Kopf und Fußzeile:
\usepackage{fancyhdr}
\pagestyle{plain}

\fancypagestyle{plain}{
	\fancyhf{}
	%Kopfzeile links
	\fancyhead[L]{\includegraphics[height = 35px]{img/Schimatta_LOGO.pdf}}
	%Kopfzeile rechts
	\fancyhead[R]{}
	%Fußzeile anzeigen
	%Fußzeile Rechts
	\fancyfoot[R]{\small{Stefan Strobel}}
	%Fußzeile Links
	\fancyfoot[C]{Seite \thepage}
	\fancyfoot[L]{Elektronik Fehlersuche}
}
% Heading for Titlepage
\fancypagestyle{myheadings}{
	\fancyhf{}
	%Kopfzeile links
	\fancyhead[L]{\includegraphics[height = 35px]{img/Schimatta_LOGO.pdf}}
	%Kopfzeile rechts
	\fancyhead[R]{}
	%Fußzeile anzeigen
	%Fußzeile Rechts
	%\fancyfoot[R]{\small{Die Drei von der Tanke}}
	%Fußzeile Links
	%\fancyfoot[L]{MCT Tutorium}
}
%Linie oben
%\renewcommand{\headrulewidth}{0.5pt}

\setlength{\headheight}{40pt}

\newlength\tindent
\setlength{\tindent}{\parindent}
\setlength{\parindent}{0pt}
\renewcommand{\indent}{\hspace*{\tindent}}

%opening
\date{\today}
\title{Fehlersuche in Elektronischen Schaltungen}
\subtitle{Grundlagen zum Finden und beheben von Problemen in moderner Elektronik}
\author{Stefan Strobel}


\begin{document}
	
	\maketitle
	
	\pagebreak
	
	\pagestyle{plain}
	
	\tableofcontents
	
	\pagebreak
	
	\input Einleitung
	
	\section{Das richtige Arbeitsumfeld}

Im Folgenden Kapitel werden Werkzeuge und Methoden aufgelistet, die zum professionellen Umgang mit elektrischen Schaltungen vorhanden bzw. angewendet werden sollten.
Über die nötige Qualität lässt sich streiten, allerdings sollte man nicht am falschen Ende sparen. Hochwertigere Geräte lassen sich für gewöhnlich auch mit weniger Erfahrung sicher bedienen.

Es gibt viele Werkzeuge und Messgeräte, die ab einem gewissen Aufwand der Reparatur benötigt werden. Allerdings sind diese häufig recht hochpreisig und zur Beherrschung ist eine umfassende Ausbildung nötig. Daher werden diese hier vorerst nicht näher beschrieben.

\subsection{Grundausstattung}

\begin{itemize}
	\item Multimeter (Unverzichtbares Hilfsmittel zum Messen verschiedenster elektrischer Größen)
	\item Lötkolben/Lötstation
	\item Heißluftlötstation (zum Tauschen von SMD Bausteinen unerlässlich)
	\item Pinzetten (feine, präzise gefertigte zur sicheren Platzierung von Bauteilen)
	\item Auflicht Mikroskop (günstige USB Varianten erhältlich)
	\item Lötzubehör
	\begin{itemize}
		\item Lötzinn (mehrere Stärken (0,5mm und 1mm mindestens))
		\item Flussmittel
		\item Lötbesteck
		\item Entlötlitze in verschiedenen Stärken
	\end{itemize}
	\item ESD Schutzausrüstung
\end{itemize}

\subsection{Arbeitsschutz}

Um sicher an elektronischen Schaltungen arbeiten zu können, sich zum einen nicht selbst zu gefährden, zum anderen nicht unbeabsichtigt weitere Fehler zu verursachen, sollten einige Grundregeln beachtet werden:

\subsubsection{Arbeiten Sie nicht unter Spannung}

Wenn es sich vermeiden lässt sollte man sämtliche Arbeiten (Diagnosen, Lötarbeiten) mit getrennter Spannungsversorgung durchführen.
Dies gilt vor allem bei Geräten, die direkt mit Netzspannung betrieben werden.
Ein Fehler hier kann ggf. tödlich enden.

Viele Elektronische Geräte arbeiten glücklicherweise mit einer sehr kleinen Spannung von idR. max 12V. Bei solch geringen Spannungen ist zwar die Verletzungsgefahr sehr gering, jedoch kann bei Arbeiten unter Spannung das Gerät, das man eigentlich reparieren will noch stärker beschädigt werden.

Verursacht man beispielsweise bei Messungen versehentlich einen Kurzschluss, kann dies bisher intakte Teile der Schaltung beschädigen.

Daher sollte soweit möglich nur Spannungsfrei gearbeitet werden.
Ist dies bei einigen Messungen nicht möglich, sollte jedoch besondere sorgfältig gearbeitet werden. Alle nicht benötigten Teile der Schaltung können beispielsweise abgedeckt werden um das Risiko von versehentlichen Berührungen zu minimieren.

Auch sollten Unterlagen beim arbeiten gewählt werden, die nichtleitend sind.
Ansonsten kann es auch hier zu ungewollten Kurzschlüssen kommen.

\subsubsection{ESD Schutz}

ESD (electrostatic discharge, oder elektrostatische Entladung) beschreibt ein häufiges Problem bei der Reparatur und Fertigung von Elektronik. Durch die immer kleiner werdenden Strukturen auf Computerchips (Speichern, Prozessoren, etc.) werden diese immer empfindlicher für kleine elektrische Entladungen. Diese treten immer dann auf, wenn sich zwei Körper unterschiedlich elektrisch aufladen. Läuft man beispielsweise mit Gummisohlen auf Teppichboden, kann es vorkommen, dass man sich beim Gehen elektrostatisch auflädt. Berührt man nun etwas leitendes, findet ein Ladungsausgleich in Form eines kleinen Lichtbogens statt (Blitz).
Trifft ein solcher Lichtbogen bei der Arbeit auf ein empfindliches elektronisches Bauteil, treten dort kurzzeitig relativ hohe Ströme auf, die das Bauteil beschädigen können.
Dies kann von kleinen Fehlfunktionen bis zu einem Totalausfall vieles zur Folge haben.
Um das zu vermeiden sollte man daher auf einen ESD gerechten Arbeitsplatz achten.

Am einfachsten kann man dies durch eine spezielle ESD Matte erreichen, die man als Arbeitsunterlage verwendet. Diese saugt praktisch alle ungewollten elektrischen Ladungen ab.
In Verbindung dazu hilft ein ESD Schutz Armband.
Dieses kann man idR. an der Matte anschließen und erdet sich dadurch selbst.
Auf diesem Weg werden auch ungewollte Ladungen vom eigenen Körper abgeleitet und so eine elektrische Entladung verhindert.
	
	\section{Fehlersuche}

\subsection{Ursachenforschung}

In der Elektronik ist häufig nicht leicht feststellbar, wo in einer komplexen Schaltung ein Fehler versteckt ist.
In vielen Fällen ist es daher hilfreich herauszufinden, was den Fehler ausgelöst haben könnte.
Soll heißen, durch welche äußeren Umstände könnte der Fehler erzeugt worden sein.

Im folgenden ein paar klassische Fehlerursachen und Probleme, die dadurch auftreten können:

\begin{itemize}
	\item Wasserschaden ==> Korrosion, Kurzschluss(Zerstörung von Bauteilen)
	\item Übertemperatur ==> Zerstörung von Bauteilen mit hoher Verlustleistung, Kontaktprobleme durch Verzug
	\item Mechanische Einwirkung ==> Bruch von Leiterbahnen, Brüche in Bauteilen, Kontaktprobleme
	\item Überspannung (Blitzeinschlag, Netzteildefekt) ==> Zerstörung von Schutzeinrichtungen, oder anderen Bauteilen
\end{itemize}



\subsection{optische Fehlersuche}

Die erste Analyse eines Fehlers sollte grundsätzlich durch eine genaue optische Fehlersuche erfolgen.
Auch wenn viele elektronische Fehler optisch nicht zu erkennen sind, führen einige Probleme doch zu sichtbaren Phänomenen.

Da elektrische Bauteile immer kleiner werden empfiehlt sich hier die Verwendung eines, wenn auch einfachen Mikroskops. Günstige Geräte mit USB Anschluss sind für schmales Geld erhältlich und erfüllen hier sehr gut Ihren Zweck.

Das genaue Augenmerk der optischen Fehlersuche hängt an dieser Stelle bereits von der vermuteten Ursache des Fehlers ab.

Auffälligkeiten:
\begin{itemize}
	\item Korrosion (sämtliche Bauteilanschlüsse absuchen, sofern zugänglich)
	\item Rückstände von Flüssigkeiten
	\item Brüche, Rissen der Leiterplatte, oder in Bauteilen
	\item Verfärbung von Teilen der Leiterplatte, oder von bestimmten Bauteilen
	\item sichtbare Beschädigungen an Bauteilen (häufig sichtbar durch kleine Löcher im Chipgehäuse) grundsätzlich sollte man sich jedes Bauteil, das in irgend einer Form merkwürdig aussieht näher untersuchen.
\end{itemize}

\subsection{messtechnische Fehlersuche}

Spannungen verfolgen

Kurzschlüsse messen

Dioden messen

Kapazitäten nachmessen (funktioniert selten gut in eingebautem Zustand)

\subsection{häufige Fehlerquellen}

Elkos

Schutzeinrichtungen(Sicherungen, Schutzdioden)

Mechanische/Kontaktprobleme (abgerissene Steckverbinder, Korrosion)

\subsection{Achtung Folgefehler}

Wenn man ein Problem löst, Vorsicht, evtl. liegt das eigentliche Problem erst dahinter...
	
	
\end{document}